% =========================================================
% Algebra of Functions — Notes
% Chapter: functions
% Section: algebra-of-functions
% Volume III · Analysis
% =========================================================

\section{Algebra of Functions}
\label{sec:algebra-of-functions}

% ---------------------------------------------------------
% Toolkit
% ---------------------------------------------------------

\begin{tcolorbox}[title={Toolkit --- Algebra of Functions}]
\begin{itemize}
  \item All results here are set-theoretic: no limits, no topology.
  \item Composition of injections is injective; composition of
    surjections is surjective.
  \item The inverse of a bijection is a bijection.
  \item Preimage distributes over $\cup$ and $\cap$ exactly:
    $f^{-1}(A \cup B) = f^{-1}(A) \cup f^{-1}(B)$ and
    $f^{-1}(A \cap B) = f^{-1}(A) \cap f^{-1}(B)$.
  \item Image distributes over $\cup$ but only $\subseteq$ over $\cap$:
    $f(A \cup B) = f(A) \cup f(B)$ but $f(A \cap B) \subseteq
    f(A) \cap f(B)$ in general.
\end{itemize}
\end{tcolorbox}

% ---------------------------------------------------------
% Reference — Intervals
% ---------------------------------------------------------

\begin{remark*}[Intervals on $\mathbb{R}$]
The standard interval types are defined in Volume II. The facts
used here: every open interval $(a,b)$ consists entirely of cluster
points; $[a,b]$ is closed and bounded (the setting for the Extreme
Value Theorem and Heine--Cantor); the natural domain of a
real-valued function is typically an interval or a union of
intervals.
\end{remark*}

% ---------------------------------------------------------
% Proposition — Composition of injections
% ---------------------------------------------------------

\begin{proposition}[Composition of Injections Is Injective]
\label{prop:composition-injective}
Let $f : A \to B$ and $g : B \to C$. If $f$ and $g$ are both
injective, then $g \circ f : A \to C$ is injective.

\smallskip
\noindent
\hyperref[prf:composition-injective]{\textit{Go to proof.}}
\end{proposition}

\begin{remark*}[Standard quantified statement]
\[
  \bigl(f \text{ injective}\bigr) \wedge \bigl(g \text{ injective}\bigr)
  \;\Rightarrow\;
  \forall a_1, a_2 \in A :\;
  g(f(a_1)) = g(f(a_2)) \Rightarrow a_1 = a_2.
\]
\end{remark*}

\begin{remark*}[Interpretation]
The proof is a direct unwinding: $g(f(a_1)) = g(f(a_2))$ implies
$f(a_1) = f(a_2)$ by injectivity of $g$, which implies $a_1 = a_2$
by injectivity of $f$.
\end{remark*}

% ---------------------------------------------------------
% Proposition — Composition of surjections
% ---------------------------------------------------------

\begin{proposition}[Composition of Surjections Is Surjective]
\label{prop:composition-surjective}
Let $f : A \to B$ and $g : B \to C$. If $f$ and $g$ are both
surjective, then $g \circ f : A \to C$ is surjective.

\smallskip
\noindent
\hyperref[prf:composition-surjective]{\textit{Go to proof.}}
\end{proposition}

\begin{remark*}[Standard quantified statement]
\[
  \bigl(f \text{ surjective}\bigr) \wedge \bigl(g \text{ surjective}\bigr)
  \;\Rightarrow\;
  \forall c \in C,\; \exists a \in A :\; g(f(a)) = c.
\]
\end{remark*}

% ---------------------------------------------------------
% Corollary — Composition of bijections
% ---------------------------------------------------------

\begin{corollary}[Composition of Bijections Is Bijective]
\label{cor:composition-bijective}
If $f : A \to B$ and $g : B \to C$ are both bijective, then
$g \circ f : A \to C$ is bijective.
\end{corollary}

% ---------------------------------------------------------
% Proposition — Inverse of bijection
% ---------------------------------------------------------

\begin{proposition}[Inverse of a Bijection Is a Bijection]
\label{prop:inverse-bijection}
Let $f : A \to B$ be a bijection. Then $f^{-1} : B \to A$ is a
bijection.

\smallskip
\noindent
\hyperref[prf:inverse-bijection]{\textit{Go to proof.}}
\end{proposition}

\begin{remark*}[Standard quantified statement]
\[
  f \text{ bijective}
  \;\Rightarrow\;
  f^{-1} \text{ bijective},\quad
  f^{-1} \circ f = \mathrm{id}_A,\quad
  f \circ f^{-1} = \mathrm{id}_B.
\]
\end{remark*}

% ---------------------------------------------------------
% Proposition — Preimage distributes over union and intersection
% ---------------------------------------------------------

\begin{proposition}[Preimage Distributes over Union and Intersection]
\label{prop:preimage-union-intersection}
Let $f : A \to B$ and let $S, T \subseteq B$. Then:
\begin{enumerate}[label=(\roman*)]
  \item $f^{-1}(S \cup T) = f^{-1}(S) \cup f^{-1}(T)$.
  \item $f^{-1}(S \cap T) = f^{-1}(S) \cap f^{-1}(T)$.
  \item $f^{-1}(S^c) = (f^{-1}(S))^c$.
\end{enumerate}

\smallskip
\noindent
\hyperref[prf:preimage-union-intersection]{\textit{Go to proof.}}
\end{proposition}

\begin{remark*}[Standard quantified statement]
\[
\begin{aligned}
(i) &\quad x \in f^{-1}(S \cup T)
     \;\iff\; f(x) \in S \cup T
     \;\iff\; f(x) \in S \vee f(x) \in T
     \;\iff\; x \in f^{-1}(S) \cup f^{-1}(T). \\
(ii) &\quad x \in f^{-1}(S \cap T)
     \;\iff\; f(x) \in S \wedge f(x) \in T
     \;\iff\; x \in f^{-1}(S) \cap f^{-1}(T).
\end{aligned}
\]
\end{remark*}

\begin{remark*}[Interpretation]
Preimage is a complete lattice homomorphism: it preserves all set
operations including complement. This contrasts sharply with the
image: $f(A \cap B) \subseteq f(A) \cap f(B)$ with strict
containment possible when $f$ is not injective. The preimage
identities are used constantly in topology and measure theory to
verify that preimages of open or measurable sets are open or
measurable.
\end{remark*}
